Early explorations \citep{wang2022git, li2022blip,li2019visualbert, sun2019videobert} of vision and language pre-trained models tend to use Bert-based \citep{liu2019roberta, koroteev2021bert} models as the language decoder. Inspired by the recent success of large language models \citep{touvron2023llama, gilardi2023chatgpt, zhao2023survey} and datasets \citep{lin2014microsoft, changpinyo2021conceptual, krishna2017visual, liu2020visual, sharma2018conceptual, srinivasan2021wit, liu2023documentclip}, many studies \citep{alayrac2022flamingo, li2023blip, li2023otter} have been focused on improving vision-language pre-trained models by integrating powerful LLMs with in-context or few-shot learning capability. More recently, some visual instruction-tuned LMMs \citep{zhu2023minigpt, liu2023visual, gong2023multimodal, dai2023instructblip} have emerged, showing excellent generalization performance in unseen VL tasks. Specifically, LLaVA \citep{liu2023visual} projects the output of a visual encoder as input to LLaMA \citep{touvron2023llama} and trains both the alignment network and the LLM on synthetic data. MiniGPT4 \citep{zhu2023minigpt} is built on BLIP-2 but uses Vicuna \citep{chiang2023vicuna} as the language decoder. It only finetunes the cross-modal alignment network on longer image captions from ChatGPT. The research approaches \citep{gong2023multimodal, dai2023instructblip} are instruction-tuned on a collection of VL datasets, but InstructBLIP \citep{dai2023instructblip} uses BLIP2 \citep{li2023blip} as the backbone while \citep{gong2023multimodal} is initialized from Flamingo \citep{alayrac2022flamingo}. mPLUG-owl \citep{ye2023mplug} finetunes LLaMA \citep{touvron2023llama} model using both text instruction data and vision-language instruction data from LLaVA \citep{liu2023visual}. In comparison, we propose a large and diverse visual instruction dataset with 16 vision and language tasks and negative instructions in different semantic levels and styles. This can help improve the robustness of current LMMs.

Although LMMs are powerful in solving VL tasks, they also suffer from the hallucination  inherited from LLM. Popular image captioning metrics like CIDEr \citep{vedantam2015cider} SPICE \citep{anderson2016spice} do not appropriately penalize hallucination. CHAIR, \citep{rohrbach2018object}, is unstable and needs complex human-crafted parsing rules for exact matching. Alternatively, \citep{li2023evaluating} converts the hallucination into a binary classification problem. However, it requires the input questions to follow specific templates, such as \textit{"Is there a/an <object> in the image?"}. In comparison, our proposed GAVIE  can evaluate model hallucination in an  open-ended manner without needing human-annotated groundtruth answers.


